% =========================================================
% PART II -- DESCRIPTIVE STATISTICS
% =========================================================

\section{Descriptive Statistics}

% ---------------------------------------------------------
% 8. MEASURES OF DISPERSION
% ---------------------------------------------------------

\begin{frame}{8. Measures of Dispersion}

\textbf{Definition}

Measures of dispersion describe the \textbf{spread or variability}
of observations around a central value.

\vspace{0.4cm}

Two datasets may have the same mean but different levels of
variability.

\vspace{0.4cm}

\textbf{Common measures}

\begin{itemize}
    \item \textbf{Range}
    \item \textbf{Variance}
    \item \textbf{Standard deviation}
    \item \textbf{Interquartile range (IQR)}
\end{itemize}

\vspace{0.3cm}

\textbf{Purpose}

They help determine how widely the observations are distributed
and how consistent the data is.

\end{frame}


% ---------------------------------------------------------
% 8.1 RANGE -- DEFINITION AND FORMULA
% ---------------------------------------------------------

\begin{frame}{8.1 Range}

\textbf{Definition}

The range is the difference between the \textbf{maximum} and
\textbf{minimum} values in a dataset.

\vspace{0.5cm}

\textbf{Formula}

\[
\boxed{
\text{Range} = x_{\max} - x_{\min}
}
\]

\vspace{0.5cm}

\textbf{Key idea}

\begin{itemize}
    \item Uses only the smallest and largest observations.
    \item Provides a simple measure of the total spread.
    \item Highly sensitive to extreme values.
\end{itemize}

\vspace{0.3cm}

\textbf{Application:}

Useful for obtaining a quick estimate of the spread of
numerical data.

\end{frame}


% ---------------------------------------------------------
% 8.1 RANGE -- EXAMPLE
% ---------------------------------------------------------

\begin{frame}{8.1 Range -- Example}

\textbf{Example}

Consider the dataset:

\[
X=\{10,20,30,40,50\}
\]

The minimum value is
\[
x_{\min}=10
\]

and the maximum value is
\[
x_{\max}=50
\]

\vspace{0.1cm}

Therefore,
\[
\text{Range}
=
50-10
=
\boxed{40}
\]

\vspace{0.1cm}

\textbf{Interpretation}

The observations span a total interval of $40$ units.

\end{frame}


% ---------------------------------------------------------
% 8.2 VARIANCE -- DEFINITION AND FORMULA
% ---------------------------------------------------------

\begin{frame}{8.2 Variance}

\textbf{Definition}

Variance measures the average squared deviation of observations
from their mean.

\vspace{0.1cm}

For a dataset $x_1,x_2,\ldots,x_n$ with mean $\bar{x}$,

\textbf{Population variance:}
\[
\boxed{
\sigma^2
=
\frac{1}{N}
\sum_{i=1}^{N}(x_i-\mu)^2
}
\]

\textbf{Sample variance:}
\[
\boxed{
s^2
=
\frac{1}{n-1}
\sum_{i=1}^{n}(x_i-\bar{x})^2
}
\]

\textbf{Key idea}

The deviations are squared so that positive and negative
deviations do not cancel each other.

\end{frame}


% ---------------------------------------------------------
% 8.2 VARIANCE -- EXAMPLE
% ---------------------------------------------------------

\begin{frame}{8.2 Variance -- Example}

\textbf{Example}

Consider:
\[
X=\{2,4,6\}
\]

First calculate the mean:
\[
\bar{x}
=
\frac{2+4+6}{3}
=
4
\]

Squared deviations:
\[
(2-4)^2=4,\qquad
(4-4)^2=0,\qquad
(6-4)^2=4
\]

Population variance:
\[
\sigma^2
=
\frac{4+0+4}{3}
=
\boxed{\frac{8}{3}}
\]

\textbf{Interpretation}

A larger variance indicates greater spread around the mean.

\end{frame}


% ---------------------------------------------------------
% 8.3 STANDARD DEVIATION -- DEFINITION AND FORMULA
% ---------------------------------------------------------

\begin{frame}{8.3 Standard Deviation}

\textbf{Definition}

Standard deviation is the positive square root of the variance.

\textbf{Population standard deviation}
\[
\boxed{
\sigma
=
\sqrt{
\frac{1}{N}
\sum_{i=1}^{N}(x_i-\mu)^2
}
}
\]
\textbf{Sample standard deviation}
\[
\boxed{
s
=
\sqrt{
\frac{1}{n-1}
\sum_{i=1}^{n}(x_i-\bar{x})^2
}
}
\]

Standard deviation is expressed in the \textbf{same units as the
original data}, unlike variance.

\end{frame}


% ---------------------------------------------------------
% 8.3 STANDARD DEVIATION -- EXAMPLE
% ---------------------------------------------------------

\begin{frame}{8.3 Standard Deviation -- Example}

\textbf{Example}

For the dataset
\[
X=\{2,4,6\},
\]
\[
\sigma^2=\frac{8}{3}
\]
Therefore,
\[
\sigma
=
\sqrt{\frac{8}{3}}
\approx
\boxed{1.63}
\]
\textbf{Interpretation}

The standard deviation indicates the typical magnitude of
deviation of observations from the mean.

\textbf{Application}

Standard deviation is widely used in data analysis and machine
learning to measure feature variability.

\end{frame}


% ---------------------------------------------------------
% 8.4 INTERQUARTILE RANGE -- DEFINITION AND FORMULA
% ---------------------------------------------------------

\begin{frame}{8.4 Interquartile Range (IQR)}

\textbf{Definition}

The interquartile range measures the spread of the
\textbf{middle 50\%} of the observations.

\vspace{0.4cm}

The data is divided into four parts using quartiles:

\begin{itemize}
    \item $Q_1$ -- First quartile (25th percentile)
    \item $Q_2$ -- Second quartile (50th percentile / median)
    \item $Q_3$ -- Third quartile (75th percentile)
\end{itemize}

\vspace{0.3cm}

\textbf{Formula}

\[
\boxed{
IQR=Q_3-Q_1
}
\]

\vspace{0.3cm}

\textbf{Key idea}

IQR focuses on the middle portion of the data and is
less affected by extreme observations.

\end{frame}


% ---------------------------------------------------------
% 8.4 IQR -- EXAMPLE
% ---------------------------------------------------------

% ---------------------------------------------------------
% 8.4 IQR -- EXAMPLE 1
% ---------------------------------------------------------

% ---------------------------------------------------------
% 8.4 IQR -- EXAMPLE 1
% ---------------------------------------------------------

\begin{frame}{8.4 Interquartile Range -- Example}

\textbf{Example}

Consider the ordered dataset:
\[
X=\{1,2,3,4,5,6,7,8,9\}
\]

The median is
\[
Q_2=5
\]

\textbf{Lower half}
\[
\{1,2,3,4\}
\]

Therefore,
\[
Q_1
=
\frac{2+3}{2}
=
\boxed{2.5}
\]

\end{frame}


% ---------------------------------------------------------
% 8.4 IQR -- EXAMPLE 2
% ---------------------------------------------------------

\begin{frame}{8.4 Interquartile Range -- Example}

\textbf{Upper half}

\[
\{6,7,8,9\}
\]

Therefore,

\[
Q_3
=
\frac{7+8}{2}
=
\boxed{7.5}
\]

From the previous example:

\[
Q_1=2.5,
\qquad
Q_3=7.5
\]

\textbf{Quartiles}

\[
Q_1=2.5,
\qquad
Q_2=5,
\qquad
Q_3=7.5
\]

\end{frame}


% ---------------------------------------------------------
% 8.4 IQR -- CALCULATION
% ---------------------------------------------------------

\begin{frame}{8.4 Interquartile Range -- Calculation}

\textbf{Formula}
\[
\boxed{IQR=Q_3-Q_1}
\]

\textbf{Calculation}

Using
\[
Q_1=2.5,
\qquad
Q_3=7.5
\]
we get
\[
IQR=7.5-2.5
\]
\[
\boxed{IQR=5}
\]

\textbf{Interpretation}

The interquartile range is $5$, meaning that the middle
$50\%$ of the observations are spread over an interval of
$5$ units.

\textbf{Key Point}

IQR is less affected by extreme values than the range.

\end{frame}

% ---------------------------------------------------------
% 8.5 COMPARISON OF DISPERSION MEASURES
% ---------------------------------------------------------

\begin{frame}{8.5 Comparison of Dispersion Measures}

\small

\begin{table}
\centering

\begin{tabular}{>{\bfseries}l l l}
\toprule
Measure & What it describes & Outlier sensitivity \\
\midrule
Range & Total spread & High \\
Variance & Squared deviation from mean & High \\
Standard deviation & Spread around mean & High \\
IQR & Spread of middle 50\% & Low \\
\bottomrule
\end{tabular}

\end{table}

\vspace{0.4cm}

\textbf{When to use?}

\begin{itemize}
    \item \textbf{Range:} Quick measure of total spread.
    \item \textbf{Variance:} Useful in statistical and mathematical models.
    \item \textbf{Standard deviation:} Easy to interpret in the original units.
    \item \textbf{IQR:} Useful for skewed data and datasets with outliers.
\end{itemize}

\end{frame}